\id{IRSTI 06.73.15}{}

{\bfseries Tax regulation of the household sector in Kazakhstan: an
analytical assessment of the resource base and tax burden}

{\bfseries \tsp{1}G.A.Maulenberdieva}\fig{e/image7}[]{\bfseries \envelope ,
\tsp{1}A.U.Abishоva}\fig{e/image8}[]
{\bfseries ,
\tsp{2}D.B.Balabekova}\fig{e/image1}[],

{\bfseries \tsp{2}G.A.
Zhadigerova}\fig{e/image8}[]{\bfseries ,\tsp{1}A.Zh.Tuleeva}\fig{e/image8}[]

\emph{\tsp{1}M. Auyezov South Kazakhstan Research
University, Shymkent, Kazakhstan,}

\emph{\tsp{2}Central Asian Innovation University, Shymkent,
Kazakhstan}

\envelope Corresponding author:
\href{file:///C:/Users/HP/Desktop/gulgar2018g@mail.ru}{gulgar2018g@mail.ru}

The article provides a comprehensive analysis of tax regulation of the
household sector in the Republic of Kazakhstan, taking into account the
current tax system, regional resource base, and income structure of the
population. Based on statistical data for 2019--2024, an assessment of
direct and indirect taxes that form the tax burden on households is
carried out, broken down by region and taxpayer category. The author
proposes an approach to calculating the absolute and relative tax burden
based on the division of taxes by source of payment - income and
consumption. The results of the study indicate significant regional
differences in the tax burden and highlight the need to improve
monitoring and the methodological basis for regulating the household
sector. The study revealed significant differences between regions in
terms of the level of tax burden on the population. It was found that
indirect taxes account for a high share of household expenditures, which
exacerbates interregional imbalances. Insufficient detail in data on
consumer structure limits the ability to accurately model the tax
burden. As a result, the calculated tax burden indicators are
conditional and dependent on the assumptions made. This underscores the
need to further improve the assessment methodology and the system for
collecting statistical information.

{\bfseries Keywords:} tax regulation, household income, tax burden, direct
taxes, indirect taxes, tax residents.

{\bfseries Қазақстандағы үй шаруашылығы секторының салық реттеуі: ресурстық
база мен салық жүктемесін аналитикалық бағалау}

{\bfseries \tsp{1}Г.А. Мауленбердиева\envelope ,
\tsp{1}А.У.Абишова, \tsp{2}Д.Б.Балабекова,
\tsp{2}Г.А. Жадигерова, \tsp{1}А.Ж.Тулеева}

\emph{\tsp{1}М. Әуезов ат. Оңтүстік Қазақстан зерттеу
университеті, Шымкент, Қазахстан,}

\emph{\tsp{2}Орталық Азия Инновациялық университеті,
Шымкент, Қазахстан,}

e-mail:
\href{file:///C:/Users/HP/Desktop/gulgar2018g@mail.ru}{gulgar2018g@mail.ru}

Мақалада қазіргі салық жүйесін, өңірлік ресурстық базаны және халықтың
табыс құрылымын ескере отырып, Қазақстан Республикасындағы үй
шаруашылығы саласын салықтық реттеуге жан-жақты талдау жасалған.
2019-2024 жылдардағы статистикалық деректер негізінде үй
шаруашылықтарына салық жүктемесін қалыптастыратын тікелей және жанама
салықтарды өңірлер мен салық төлеушілердің санаттары бөлінісінде бағалау
орындалды. Салықтарды төлем көздері - кірістер мен тұтыну бойынша бөлуге
негізделген абсолютті және салыстырмалы салық жүктемесін есептеуге
авторлық тәсіл ұсынылды. Зерттеу нәтижелері салық жүктемесіндегі
айтарлықтай аймақтық айырмашылықтарды көрсетеді және үй шаруашылығы
секторын реттеудің мониторингі мен әдіснамалық базасын жетілдіру
қажеттілігін көрсетеді. Зерттеу халыққа салынатын салық жүктемесінің
деңгейі бойынша аймақтар арасындағы айтарлықтай айырмашылықтарды
анықтады. Үй шаруашылықтарының шығындарында жанама салықтар жоғары
үлесті алатыны анықталды, бұл аймақаралық теңгерімсіздіктерді күшейтеді.
Тұтынушы құрылымы туралы мәліметтердің жеткіліксіз егжей-тегжейлері
салық ауыртпалығын дәл модельдеу мүмкіндіктерін шектейді. Нәтижесінде
салық жүктемесінің есептік көрсеткіштері шартты және қабылданған
болжамдарға тәуелді болады. Бұл статистикалық ақпаратты жинау жүйесін
бағалау және жақсарту әдістемесін одан әрі жетілдіру қажеттілігін
көрсетеді.

{\bfseries Түйін сөздер}: салықтық реттеу, үй шаруашылықтарының табысы,
салық жүктемесі, тікелей салықтар, жанама салықтар, салық резиденттері.

{\bfseries Налоговое регулирование сектора домашних хозяйств в Казахстане:
аналитическая оценка ресурсной базы и налоговой нагрузки}

{\bfseries \tsp{1}Г.А.Мауленбердиева\envelope ,
\tsp{1}А.У.Абишова, \tsp{2}Д.Б.Балабекова,}

{\bfseries \tsp{2}Г.А.Жадигерова,
\tsp{1}А.Ж.Тулеева}

\emph{\tsp{1}Южно-Казахстанский исследовательский
университет им. М. Ауезова, Шымкент, Казахстан,}

\emph{\tsp{2}Центрально-Азиатский инновационный университет,
Шымкент, Казахстан,}

e-mail:
\href{file:///C:/Users/HP/Desktop/gulgar2018g@mail.ru}{gulgar2018g@mail.ru}

В статье проведен комплексный анализ налогового регулирования сектора
домашних хозяйств в Республике Казахстан с учетом действующей налоговой
системы, региональной ресурсной базы и структуры доходов населения. На
основе статистических данных 2019--2024 гг. выполнена оценка прямых и
косвенных налогов, формирующих налоговую нагрузку на домохозяйства, в
разрезе регионов и категорий налогоплательщиков. Предложен авторский
подход к расчету абсолютной и относительной налоговой нагрузки,
основанный на разделении налогов по источникам уплаты - доходам и
потреблению. Результаты исследования указывают на существенные
региональные различия в налоговой нагрузке и подчеркивают необходимость
совершенствования мониторинга и методологической базы регулирования
сектора домашних хозяйств. Исследование выявило значительные различия
между регионами по уровню налоговой нагрузки на население. Было
установлено, что в расходах домашних хозяйств высокую долю занимают
косвенные налоги, что усиливает межрегиональные дисбалансы.
Недостаточная детализация данных о потребительской структуре
ограничивает возможности для точного моделирования налогового бремени. В
результате расчётные показатели налоговой нагрузки оказываются условными
и зависимыми от принятых допущений. Это подчёркивает необходимость
дальнейшего совершенствования методологии оценки и улучшения системы
сбора статистической информации.

{\bfseries Ключевые слова}: налоговое регулирование, доходы домохозяйств,
налоговая нагрузка, прямые налоги, косвенные налоги, налоговые
резиденты.

{\bfseries Introduction.} The household sector plays a key role in the
sustainable socio-economic development of Kazakhstan, as it is the main
taxpayer of both direct and indirect taxes, which form a significant
part of state and local budget revenues. In the context of increasing
uncertainty, accelerating inflationary processes and changes in the
income structure of the population, there is a growing need to develop
and improve tax regulation mechanisms that ensure a balance of public
and private interests.

However, despite the importance of the sector, Kazakhstan still lacks a
methodologically sound system for assessing household participation in
tax relations, as well as a systematic database on the tax burden of the
population. In this regard, the study of the tax burden and the resource
base of the household sector is an urgent area of scientific analysis.

The relevance of the study lies in the growing role of households in
generating tax revenues, coupled with the lack of effective systems for
monitoring the actual tax base of the population. Increasing income
differentiation and the need to assess the impact of taxes on
consumption patterns require the development of a unified methodological
approach to analyzing the tax burden. This makes the study important for
improving tax policy and adapting it to socio-regional conditions. Thus,
the study aims to fill the gaps that hinder the development of an
economically sound tax policy for the household sector.

In the world of literature, the issue of household tax burden is
explored through the lens of concepts such as tax burden (Musgrave,
1980), optimal tax system (Mirrlees, 2011), consumer behavior (Becker,
1996), and income distribution (Piketty, 2014) {[}1-4{]}. Researchers
from the EU, the United States, and the OECD focus on comparing the tax
burdens of households with different income levels, examining the impact
of indirect taxes on consumption and inequality, and assessing the
effectiveness of current tax benefits and deductions {[}5-7{]}.

In Kazakhstan, theoretical foundations of taxation have been explored in
the works of A.A. Nurumov, M.K. Kozhakhmetova and M.S. Dosmanbetova.
However, there has been no comprehensive assessment of household tax
burdens. Additionally, there are no studies that compare regional
resource bases and tax revenues while considering socio-demographic
factors {[}8-10{]}.

Modern global discourse on household taxation increasingly focuses on
the actual distribution of the tax incidence. Bachas et al. (2024),
writing in the Review of Economic Studies (h-index \textgreater{} 170),
highlight the complexities of consumption taxes in developing economies,
noting that the real burden is often distorted by consumption patterns
and informal sectors {[}11{]}. In the context of resource dependency and
regional development, the work of Li Hong (2019) in the Journal of
Cleaner Production (h-index \textgreater{} 200) is particularly
relevant. The author proposes an "environmental scarcity tax" model that
utilizes a compensation system for ecosystem services to balance the
fiscal load in resource-intensive regions. This approach correlates with
the disparities identified in Kazakhstan's Atyrau and Mangistau regions,
where fiscal indicators require a more nuanced interpretation than
standard direct income taxation {[}12{]}. Furthermore, research by
Benzarti and Carloni (2019) in the American Economic Journal (h-index
\textgreater{} 180) emphasizes that VAT reforms are unevenly translated
into consumer welfare, supporting our hypothesis regarding the hidden
effects of indirect taxes on the budgets of Kazakhstani citizens who
spend over 90\% of their income on goods and services{[} 13{]}.

To move beyond a descriptive analysis, this study addresses the
following research questions:

\emph{RQ1:} How does the regional specificity of the resource base and
household consumption structure affect the actual distribution of the
tax burden in Kazakhstan?

\emph{RQ2:} To what extent do indirect taxes create hidden imbalances in
the tax system?

\emph{Hypothesis:} The current system of tax monitoring in Kazakhstan
fails to fully account for the high share of indirect taxes in household
expenditures (reaching up to 92.8\%), leading to a "hidden regressivity"
of the tax system and significant regional disparities, particularly in
cities of republican significance and resource-rich regions.

The proposed work expands on existing research and is one of the first
integrated approaches to assessing the resource base of the household
sector.

The purpose of the study is to develop a methodological approach for
assessing the tax burden of households in Kazakhstan and analyzing their
resource base, taking into account regional specifics. To achieve this
goal, the mechanisms of tax regulation and the structure of existing
taxes are being studied, regional differences are being identified, a
methodology for calculating the tax burden is being developed, and
conclusions and recommendations are being formulated to improve
monitoring of the sector.

{\bfseries Materials and methods.} This study is a quantitative empirical
and analytical analysis, including a descriptive study of tax revenues
and the resource base of households, the use of economic and statistical
methods and modeling the tax burden. The work has an applied
socio-economic character and is aimed at assessing the level of the tax
burden of households in Kazakhstan and the factors influencing its
formation. The object of the study is the household sector of the
Republic of Kazakhstan, which is a collection of individuals involved in
the formation of tax revenues. These include payers of individual income
taxes, indirect taxes included in the cost of goods and services, as
well as property taxes such as taxes on property, land, and
transportation.

The work uses data from the Bureau of National Statistics of the
Republic of Kazakhstan, including indicators of population, employment,
as well as the structure of income and expenses for 2019-2024.
Additionally, information from the Ministry of Finance of the Republic
of Kazakhstan on regional income tax receipts, VAT, excise taxes and
property taxes for 2019-2024 has been applied. The data of the Agency
for Strategic Planning and Reforms on the distribution of the population
by income groups and regulatory rates from the Tax Code of the Republic
of Kazakhstan were also used. The information was selected using a
targeted and continuous method, covering all regions of Kazakhstan,
including Astana, Almaty and Shymkent.

The analysis includes complete data on tax revenues according to the
relevant budget classification codes. In addition, updated data on the
structure of income and expenses of the population for the period
2019-2024 have been taken into account. The study uses a set of data
processing methods, including statistical analysis of the dynamics of
indicators, inter-regional comparisons and calculation of averages. The
data was grouped by type of taxes, regions, and income levels of the
population. Economic and mathematical models were also used to calculate
the absolute and relative tax burden, supplemented by the
author' s models for estimating indirect taxes. The
results were visualized using graphs and diagrams, which made it
possible to conduct a comparative analysis of regional differences and
peculiarities of taxation of residents and non-residents.

The choice of methods is explained by the purpose of the study aimed at
developing a model for assessing the tax burden of households. The
aggregated statistical data used makes it possible to apply economic,
statistical and analytical approaches. In addition, the specifics of
Kazakhstan' s tax system, which includes both direct and
indirect taxes, require the use of comprehensive analysis methods. In
the absence of detailed statistics on consumption by product groups, the
method of averaged calculated rates was chosen, which corresponds to
accepted international practices for estimating indirect taxes (OECD
Taxation Consumption Framework) {[}14{]}.

In a narrow sense, tax regulation of the household sector is a
purposeful government influence on the financial resources of households
through the establishment and collection of taxes, supported by the
institution of state coercion, in order to implement fiscal, regulatory,
social and control functions. According to the authors, its purpose is
twofold: to promote the strategic objectives of the
country' s development and to ensure a balance between
the private and public interests of taxpayers and consumers of public
services. Tax regulation is carried out through a special mechanism that
includes subjects and objects of taxation, forms of interaction, methods
and instruments of influence (Figure 1). This mechanism ensures the
coordination and harmonization of the functions of the tax system within
the framework of the country' s development priorities.

{\bfseries Fig.1 - The mechanism of tax regulation of the household
sector}

\emph{Note: compiled by the authors}

An analysis of the current legislation on taxes and fees has shown that
currently individuals are payers of taxes (fees) on income, property and
actually payers of taxes on consumption (indirect or transferable
taxes). At the same time, households pay an individual income tax, which
reduces their disposable income, as well as indirect taxes such as VAT
and excise taxes, which increase consumer spending. In addition, taxes
on property, land, and transportation increase the cost of owning
assets.

The methodology shifts from simple statistical reporting to analytical
modeling of tax incidence. By employing an author-developed approach,
the study distinguishes between taxes by source of payment (income vs.
consumption). In the absence of detailed consumption data, the method of
averaged calculated rates was adopted, aligning with the OECD Taxation
Consumption Framework to estimate the actual burden on the population.

{\bfseries Results and discussions.} For the budget system of Kazakhstan,
the most significant taxes paid by households are personal income tax,
vehicle tax on individuals, VAT, and excise taxes. According to the
Ministry of Finance of the Republic of Kazakhstan, taxes on goods sold
in the Republic in 2024 amounted to 32.2\% of all state budget revenues,
taxes on imported goods -- 11.1\%, which respectively amounted to 7.5\%
and 2.6\% of Kazakhstan' s GDP in 2024.

In addition, 33.1\% of all revenues of local budgets of the Republic of
Kazakhstan' s constituent entities come from personal
income tax (which is even more than income tax revenues, which account
for only 21.1\% of local budget revenues), 7.5\% come from property
taxes (transport tax and property tax), and 9.7\% comes from excise
taxes on goods produced in the Republic of Kazakhstan {[}15{]}.

The dynamics of personal income tax receipts to local (regional,
district) budgets of the subjects of the Republic of Kazakhstan in
2019-2024 was stable (Figure 2).

{\bfseries Fig.2 - Dynamics of personal income tax receipts to local
(regional, district) budgets of the subjects of the Republic of
Kazakhstan in 2020-2024, thousand tenge}

\emph{Note: compiled by the authors}

As shown in Figure 2, in terms of personal income tax, the highest
revenues to local budgets in Kazakhstan' s regions come
from Almaty, Astana, and the Atyrau region. The outsiders are the three
new regions formed in 2022: Zhetysu, Abai, and Ulytau, as well as
Kyzylorda and North Kazakhstan regions. It should also be noted that
over the past six years, there has been an increase in tax revenues from
personal income tax in all regions except the Atyrau region {[}16{]}.

There are no discernible patterns in local budget revenues from personal
income tax on vehicles. Almost all regions show a clear downward trend
in vehicle tax revenues in 2020 compared to 2019, but starting in 2021,
there is a growth trend in revenues. In terms of the amount of vehicle
tax revenue from individuals, the leaders are Almaty, Almaty Region, and
Astana.

VAT revenues to the republican budget by province and region are shown
in Figure 3. Of course, it should be noted that individuals were not the
final payers of all VAT amounts. Determining the actual share of VAT
revenues from Kazakhstanis' personal budget is not possible.

{\bfseries Fig.3 - Dynamics of VAT revenues to the republican budget of
the republic' s constituent entities in 2019-2023,
thousand tenge}

\emph{Note: compiled by the authors}

There is a trend toward increased VAT revenues in the budget system
(which, given the unchanged tax rates, correlates well with the growth
in consumption by the population due to rising prices). Figure 3 shows
VAT revenues by budget classification codes from ``105101'' to
``105116'' inclusive {[}16{]}.

Excise tax revenues from excisable goods to the republican budget vary
significantly across regions, demonstrating a high concentration of
revenues in certain regions and cities of republican significance. The
highest values are observed in the Atyrau region, where revenues
significantly exceed those of other regions in all years, which is
probably due to the concentration of the oil and gas industry. High
levels of revenue are also characteristic of Almaty, Astana, and the
Mangistau region, where steady growth can be seen from year to year.

Most other regions show relatively low tax revenues, but moderate growth
can also be seen there, especially in the Karaganda, Akmola, Pavlodar,
and East Kazakhstan regions. The different trends are not related to the
consumption of excisable goods by the population, but to the territorial
location of businesses whose goods (mainly alcohol and tobacco products)
are subject to excise taxation. Nevertheless, in 2024, almost all
regions will see an increase in revenues compared to previous years,
reflecting the overall trend of growth in the tax base.

The analysis of the resource base for tax regulation of the domestic
household sector in terms of taxation is not related to potential or
actual entrepreneurial activity. The determination of the potential and
actual number of taxpayers raises certain issues, since individuals have
tax capacity (i.e., the ability to have tax rights and bear tax
obligations) from birth. Consequently, the entire population of the
country is potential taxpayers. In reality, however, those individuals
who have a taxable object -- income, property -- become taxpayers, and
the tax authorities keep records of the number of taxpayers in terms of
taxes. Thus, the potential number of taxpayers is the total population
of Kazakhstan, while the actual number varies for each direct tax.
Indirect taxes are paid by all individuals, since in a market economy it
is impossible to live without entering into (personally or through
representatives - legal and authorized) monetary and material exchange
relations, i.e., without purchasing goods, works, and services.

According to data for 2024, the population of Kazakhstan (the potential
number of taxpayers) was about 9.7 million people, of which 63\% (6.1
million people) were urban residents and 37\% (3.6 million people) were
rural residents. At the same time, 76\% of the working-age population
aged 15 and above in 2024 was employed, and 24\% of the population was
self-employed. At the end of 2024, the largest number of Kazakhstani
citizens were employed in the service sector (62.9\% of all employed
persons), the smallest number were employed in agriculture (only 9.3\%
of employed persons), and 27.8\% were employed in industry {[}17{]}.

In the most general terms, the structure of Kazakhstan' s
population by income, according to data from the Agency for Strategic
Planning and Reforms of the Republic of Kazakhstan, is characterized by
a predominance of the population with incomes of 256,000 tenge/month,
with the smallest number (4.03\%) earning up to 47,000 tenge/month
(Figure 4).

{\bfseries Fig.4 - Income distribution by 10\% population groups}

\emph{Note: compiled by the authors}

The current resource base of tax relations of individuals by region in
Kazakhstan (Table 1) is characterized by the following features:

{\bfseries Table 1 - Indicators of the resource base of tax relations of
individuals by region in Kazakhstan}

%% \begin{longtable}[]{@{}
%%   >{\raggedright\arraybackslash}p{(\linewidth - 16\tabcolsep) * \real{0.1441}}
%%   >{\centering\arraybackslash}p{(\linewidth - 16\tabcolsep) * \real{0.1034}}
%%   >{\centering\arraybackslash}p{(\linewidth - 16\tabcolsep) * \real{0.0887}}
%%   >{\centering\arraybackslash}p{(\linewidth - 16\tabcolsep) * \real{0.1032}}
%%   >{\centering\arraybackslash}p{(\linewidth - 16\tabcolsep) * \real{0.0884}}
%%   >{\centering\arraybackslash}p{(\linewidth - 16\tabcolsep) * \real{0.1181}}
%%   >{\centering\arraybackslash}p{(\linewidth - 16\tabcolsep) * \real{0.1182}}
%%   >{\centering\arraybackslash}p{(\linewidth - 16\tabcolsep) * \real{0.1034}}
%%   >{\centering\arraybackslash}p{(\linewidth - 16\tabcolsep) * \real{0.1327}}@{}}
%% \toprule\noalign{}
%% \multirow{2}{=}{\begin{minipage}[b]{\linewidth}\centering
%% Region/area
%% \end{minipage}} &
%% \multicolumn{3}{>{\centering\arraybackslash}p{(\linewidth - 16\tabcolsep) * \real{0.2952} + 4\tabcolsep}}{%
%% \begin{minipage}[b]{\linewidth}\centering
%% Population, million people
%% \end{minipage}} &
%% \multirow{2}{=}{\centering\arraybackslash \begin{minipage}[b]{\linewidth}\centering
%% {\bfseries Average per capita cash income}
%% \end{minipage}} &
%% \multicolumn{3}{>{\centering\arraybackslash}p{(\linewidth - 16\tabcolsep) * \real{0.3396} + 4\tabcolsep}}{%
%% \begin{minipage}[b]{\linewidth}\centering
%% {\bfseries To local (regional) budgets of Kazakhstan' s
%% constituent entities, thousand tenge}
%% \end{minipage}} &
%% \multirow{2}{=}{\centering\arraybackslash \begin{minipage}[b]{\linewidth}\centering
%% {\bfseries Republican budget (VAT on goods sold in Kazakhstan)}
%% \end{minipage}} \\
%% & \begin{minipage}[b]{\linewidth}\centering
%% {\bfseries Economically active}
%% \end{minipage} & \begin{minipage}[b]{\linewidth}\centering
%% {\bfseries Employed}
%% \end{minipage} & \begin{minipage}[b]{\linewidth}\centering
%% {\bfseries IE (self-employed)}
%% \end{minipage} & & \begin{minipage}[b]{\linewidth}\centering
%% {\bfseries PIT}
%% \end{minipage} & \begin{minipage}[b]{\linewidth}\centering
%% {\bfseries Transport tax on individuals}
%% \end{minipage} & \begin{minipage}[b]{\linewidth}\centering
%% {\bfseries Property tax on individuals}
%% \end{minipage} \\
%% \midrule\noalign{}
%% \endhead
%% \bottomrule\noalign{}
%% \endlastfoot
%% Total for the Republic & 9682380 & 9233828 & 2213679 & 210840 &
%% 2213223059 & 95073209 & 19859134 & 2267018144 \\
%% Abai & 306551 & 292231 & 112578 & 182647 & 43116728 & 2349620 & 115158 &
%% -25905752 \\
%% Akmola & 422745 & 403771 & 105 020 & 190666 & 63156968 & 3342234 &
%% 355603 & 44825849 \\
%% Aktobe & 483640 & 460959 & 78 576 & 184934 & 81344861 & 3524789 & 885514
%% & 38110486 \\
%% Almaty & 753824 & 718525 & 220 541 & 144226 & 106032426 & 8297866 &
%% 435900 & 128823730 \\
%% Atyrau & 369964 & 351993 & 50 731 & 336743 & 155804752 & 2646698 &
%% 691815 & 251616996 \\
%% East Kazakhstan & 385756 & 368032 & 67563 & 224625 & 72166368 & 3843257
%% & 321620 & -115908026 \\
%% Zhambyl & 572030 & 544665 & 182281 & 141853 & 53166797 & 4605435 &
%% 467273 & 38306399 \\
%% Zhetisu & 311475 & 296405 & 79074 & 141318 & 30953587 & 2984268 & 200104
%% & 25040653 \\
%% West Kazakhstan & 354286 & 337283 & 89371 & 193094 & 64597997 & 2285310
%% & 527708 & 1040821 \\
%% Karaganda & 575570 & 552089 & 67564 & 232251 & 119216623 & 5432972 &
%% 671190 & 49624438 \\
%% Kyzylorda & 350954 & 334208 & 222548 & 150358 & 40387274 & 2436325 &
%% 207793 & 65832711 \\
%% Kostanay & 463919 & 442714 & 113390 & 204853 & 72274399 & 3034383 &
%% 331428 & 33803674 \\
%% Mangistau & 367053 & 348757 & 40422 & 243627 & 106229811 & 4482791 &
%% 418627 & 91562764 \\
%% Pavlodar & 405054 & 385838 & 59461 & 223731 & 78155898 & 2954480 &
%% 396833 & 61193645 \\
%% North Kazakhstan & 288221 & 275333 & 68780 & 202941 & 37444221 & 2280126
%% & 201804 & 24922718 \\
%% Turkestan & 849992 & 809375 & 374035 & 113566 & 67599205 & 5607273 &
%% 413675 & 96043824 \\
%% Ulytau & 100687 & 96599 & 10002 & 284793 & 38601767 & 29826 & 1146236 &
%% -33261198 \\
%% Shymkent city & 468719 & 446299 & 147563 & 141625 & 71991747 & 5556809 &
%% 2627634 & 95425575 \\
%% Almaty city & 1134959 & 1082982 & 129115 & 338663 & 600959288 & 18293411
%% & 5553138 & 891542015 \\
%% Astana city & 716981 & 685770 & 579818 & 296337 & 310022345 & 9968~927 &
%% 5006491 & 504376824 \\
%% \multicolumn{9}{@{}>{\raggedright\arraybackslash}p{(\linewidth - 16\tabcolsep) * \real{1.0000} + 16\tabcolsep}@{}}{%
%% \emph{Note: the table is based on the source {[}15,16{]}}} \\
%% \end{longtable}

According to Table 1, Almaty is the leader in all indicators considered
-- total population, working-age population, employed population,
average per capita income, and income tax, vehicle tax on individuals,
and VAT revenues. The newly formed regions of Ulytau and Zhetisu are the
outsiders in all indicators.

According to data on the structure of the population by income group,
the main taxpayers in terms of personal income tax and indirect taxes
are undoubtedly families belonging to the group with an income of
256,000 tenge per month per family member (more than 25.55\% of the
total population in Kazakhstan).

The structure of Kazakhstani expenditures is characterized by a
predominance of expenditures on the purchase of goods and payment for
services (in 2024, these expenditures amounted to almost 91.8\%), while
the least amount of funds in 2023 was spent on paying taxes - only 0.2\%
of all income (Table 2).

{\bfseries Table 2- Structure of household cash expenditures in Kazakhstan
for 2019--2024, \%}

%% \begin{longtable}[]{@{}
%%   >{\centering\arraybackslash}p{(\linewidth - 14\tabcolsep) * \real{0.0897}}
%%   >{\centering\arraybackslash}p{(\linewidth - 14\tabcolsep) * \real{0.1046}}
%%   >{\centering\arraybackslash}p{(\linewidth - 14\tabcolsep) * \real{0.1346}}
%%   >{\centering\arraybackslash}p{(\linewidth - 14\tabcolsep) * \real{0.1162}}
%%   >{\centering\arraybackslash}p{(\linewidth - 14\tabcolsep) * \real{0.1942}}
%%   >{\centering\arraybackslash}p{(\linewidth - 14\tabcolsep) * \real{0.1231}}
%%   >{\centering\arraybackslash}p{(\linewidth - 14\tabcolsep) * \real{0.1196}}
%%   >{\centering\arraybackslash}p{(\linewidth - 14\tabcolsep) * \real{0.1180}}@{}}
%% \toprule\noalign{}
%% \multirow{3}{=}{\centering\arraybackslash \begin{minipage}[b]{\linewidth}\centering
%% \end{minipage}} &
%% \multirow{3}{=}{\centering\arraybackslash \begin{minipage}[b]{\linewidth}\centering
%% Cash income - total
%% \end{minipage}} &
%% \multicolumn{6}{>{\centering\arraybackslash}p{(\linewidth - 14\tabcolsep) * \real{0.8057} + 10\tabcolsep}@{}}{%
%% \begin{minipage}[b]{\linewidth}\centering
%% Including
%% \end{minipage}} \\
%% & &
%% \multirow{2}{=}{\centering\arraybackslash \begin{minipage}[b]{\linewidth}\centering
%% {\bfseries Consumer spending}
%% \end{minipage}} &
%% \multicolumn{3}{>{\centering\arraybackslash}p{(\linewidth - 14\tabcolsep) * \real{0.4335} + 4\tabcolsep}}{%
%% \begin{minipage}[b]{\linewidth}\centering
%% {\bfseries Of which}
%% \end{minipage}} &
%% \multirow{2}{=}{\centering\arraybackslash \begin{minipage}[b]{\linewidth}\centering
%% {\bfseries Tax payments}
%% \end{minipage}} &
%% \multirow{2}{=}{\centering\arraybackslash \begin{minipage}[b]{\linewidth}\centering
%% {\bfseries Other expenses}
%% \end{minipage}} \\
%% & & & \begin{minipage}[b]{\linewidth}\centering
%% {\bfseries Food products}
%% \end{minipage} & \begin{minipage}[b]{\linewidth}\centering
%% {\bfseries Non-food products}
%% \end{minipage} & \begin{minipage}[b]{\linewidth}\centering
%% {\bfseries Paid services}
%% \end{minipage} \\
%% \midrule\noalign{}
%% \endhead
%% \bottomrule\noalign{}
%% \endlastfoot
%% 2019 & 100 & 92,8 & 49,9 & 23,7 & 19,2 & 0,3 & 6,9 \\
%% 2020 & 100 & 92,8 & 54,3 & 23,8 & 14,7 & 0,3 & 6,9 \\
%% 2021 & 100 & 93,3 & 52,9 & 24,8 & 15,6 & 0,2 & 6,5 \\
%% 2022 & 100 & 92,7 & 50,8 & 25,4 & 16,5 & 0,2 & 7,1 \\
%% 2023 & 100 & 91,9 & 50,8 & 24,0 & 17,1 & 0,2 & 7,9 \\
%% 2024 & 100 & 91,8 & 50,0 & 24,0 & 17,8 & 0,2 & 8,0 \\
%% \multicolumn{8}{@{}>{\centering\arraybackslash}p{(\linewidth - 14\tabcolsep) * \real{1.0000} + 14\tabcolsep}@{}}{%
%% \emph{Note: the table is based on the source {[}17,18{]}}} \\
%% \end{longtable}

Over the past six years, the structure of household expenditures has
remained unchanged (Table 2). Kazakhstani citizens spend the most
(49-54\%) on food, with the remainder (23-26\%) spent on non-food items.

Calculations show that, on average, about 92\% of Kazakhstani
citizens'{} budgets go toward purchasing goods and paying
for services. It is these expenditures that contain the indirect tax
component. This means that it is possible to calculate how much
conditional taxes are paid by the average Kazakhstani citizen. In tax
theory, the terms ``tax burden'' and ``tax load'' are used for such
situations.

This high concentration of expenditures on goods and services (over
90\%) aligns with the findings of Bachas et al. (2024), who argue that
in developing economies, consumption taxes represent the most
significant fiscal interface between the state and households. Such a
structure implies that the population' s tax burden is
primarily driven by consumption rather than direct income, reinforcing
the importance of modeling indirect tax incidence accurately {[}11{]}.

According to the authors, the tax burden on households (consumers of
goods, works, and services) can be calculated as follows:

\emph{ATB =}\(\sum\) \emph{DT +}\(\sum\) \emph{IT,}

\emph{RTB =}\(\frac{ATB}{Income}\)

where \emph{ATB} is the absolute tax burden (taxes paid by households).
This indicator includes the amounts of direct taxes (\emph{DT}) and
indirect taxes \emph{(IT)} (the burden of which is actually borne by
individuals);

\emph{RTB} - relative tax burden - is the ratio of the absolute tax
burden to household income (or, in other words, the share of tax
payments in an individual' s income).

Based on the proposed formulas, two options for calculating the tax
burden on households can be considered.

The first option is from the perspective of the government (tax
authorities):

\(\sum\)\emph{DT} - the amount of direct taxes collected by budgets at
all levels from a given territory;

\(\sum\)\emph{IT} - the amount of indirect taxes collected by budgets at
all levels from a given territory.

The second option is from the perspective of the household:

\(\sum\) \emph{DT\tsb{1}} - the amount of direct taxes paid by
households;

\(\sum\) \emph{IT\tsb{2}} - the amount of indirect taxes paid
by households when purchasing goods, works, and services {[}6{]}.

The logic of the calculations is based on the fact that at the first
stage (when disposable income is formed), gross income is reduced by the
amount of income tax (in this case, by the amount of personal income tax
withheld by the employer, who is the tax agent). There are two possible
scenarios here: personal income tax is levied on the income of a tax
resident or on the income of a tax non-resident. In the first case, the
tax rate will be 10\%, and in the second, 20\%. We will not consider tax
deductions due to their insignificant multiplier effect.

Next comes the stage of spending disposable income, i.e., the income
remaining to the population after taxation. As mentioned earlier,
consumption currently accounts for about 90\% of all income in
Kazakhstan. Given that the consumption of goods and services by the
population is subject to indirect taxation, which increases the price of
goods, it can be assumed that the share of indirect taxes paid by the
population will be proportional to the tax rate. However, without data
on the structure of household consumption broken down by taxable groups
of goods, we must proceed from the average estimated rates of indirect
taxes included in the price of goods. Let us assume that the estimated
average VAT rate is 12\%. It is not yet possible to calculate the share
of excise taxes and import duties in the total amount of household
consumption expenditure, so we will not take them into account.

Thus, we obtain:

\def\labelenumi{\arabic{enumi})}

1. The amount of income tax for different groups of taxpayers is as
follows:

\(\sum\) \emph{DT\tsb{r} = 0,1*TI /1,1;}

\(\sum\) \emph{DT\tsb{nr} = 0,2*TI /1,2;}

where \(\sum\) \emph{DT\tsb{r}} -- amount of direct taxes for
tax residents; \emph{TI}-- taxable income; \(\sum\)
\emph{DT\tsb{nr}} -- amount of direct taxes for non-residents.

\def\labelenumi{\arabic{enumi})}
\setcounter{enumi}{1}

1. \(\sum\) \emph{IT = 0,8* I\tsb{AT} *0,12;}

\emph{I\tsb{AT} = TI -- T\tsb{i},}

where IAT is income after taxation, \emph{T\tsb{i}} -- income
tax.

Therefore, indirect tax amounts can be calculated using the following
formulas:

\(\sum\) \emph{IT\tsb{r}} = \emph{0,8(TI - 0,1*TI)/ 1,13;}

\(\sum\) \emph{IT\tsb{nr} = 0,8(TI - 0,2 *TI /1,3) *0,12}.

3) \emph{ATB} consists of:

\emph{ATB \tsb{r}=0,1*TI /1,1+0,8* I\tsb{AT}}
*\emph{0,12=0,1*TI /1,1+0,8(TI -0,1* TI /1,1)0,12;}

\emph{ATB \tsb{nr}=0,2*TI /1,2+0,8* I\tsb{AT}}
*\emph{0,12=0,2*TI /1,2+0,8(TI -0,2*TI /1,2)0,12.}

If we take the average per capita income of the population as the
taxable income, we can calculate the absolute and relative tax burden
for tax residents and non-residents of the Republic of Kazakhstan. The
results of the calculations are presented in Tables 3 and 4. As the
calculations show, the tax burden differs from the perspective of the
state and individual taxpayers. The highest tax burden is on residents
of Almaty (79.6\%), and the lowest is in the Ulytau region (3.3\%). This
indicates a significant difference in the consumption of goods and
services by the population of these regions.

The stark regional disparities in tax burden-reaching 79.6\% in Almaty
compared to 3.3\% in Ulytau -necessitate a more nuanced approach to
fiscal monitoring. This correlates with the
' environmental scarcity'{} and
resource-based tax models proposed by Li Hong (2019) {[}12{]}. Much like
the regional ecological footprints studied by Li, Kazakhstan's
industrial regions (e.g., Atyrau and Mangistau) show high indirect tax
potential that requires a specialized compensation system to balance
regional social development with fiscal extraction.

{\bfseries Table 3 - Tax burden on households (from government
perspective)}

%% \begin{longtable}[]{@{}
%%   >{\raggedright\arraybackslash}p{(\linewidth - 12\tabcolsep) * \real{0.1762}}
%%   >{\centering\arraybackslash}p{(\linewidth - 12\tabcolsep) * \real{0.1199}}
%%   >{\centering\arraybackslash}p{(\linewidth - 12\tabcolsep) * \real{0.1347}}
%%   >{\centering\arraybackslash}p{(\linewidth - 12\tabcolsep) * \real{0.1647}}
%%   >{\centering\arraybackslash}p{(\linewidth - 12\tabcolsep) * \real{0.1499}}
%%   >{\centering\arraybackslash}p{(\linewidth - 12\tabcolsep) * \real{0.1497}}
%%   >{\centering\arraybackslash}p{(\linewidth - 12\tabcolsep) * \real{0.1047}}@{}}
%% \toprule\noalign{}
%% \multirow{2}{=}{\begin{minipage}[b]{\linewidth}\raggedright
%% Regions
%% \end{minipage}} &
%% \multicolumn{2}{>{\centering\arraybackslash}p{(\linewidth - 12\tabcolsep) * \real{0.2546} + 2\tabcolsep}}{%
%% \begin{minipage}[b]{\linewidth}\centering
%% Direct taxes, thousand tenge
%% \end{minipage}} & \begin{minipage}[b]{\linewidth}\centering
%% Indirect taxes, thousand tenge
%% \end{minipage} &
%% \multirow{2}{=}{\centering\arraybackslash \begin{minipage}[b]{\linewidth}\centering
%% Remuneration, thousand tenge
%% \end{minipage}} &
%% \multirow{2}{=}{\centering\arraybackslash \begin{minipage}[b]{\linewidth}\centering
%% Absolute tax burden, thousand tenge
%% \end{minipage}} &
%% \multirow{2}{=}{\centering\arraybackslash \begin{minipage}[b]{\linewidth}\centering
%% Relative tax burden, \%
%% \end{minipage}} \\
%% & \begin{minipage}[b]{\linewidth}\centering
%% {\bfseries PIT}
%% \end{minipage} & \begin{minipage}[b]{\linewidth}\centering
%% {\bfseries Transport tax}
%% \end{minipage} & \begin{minipage}[b]{\linewidth}\centering
%% {\bfseries VAT}
%% \end{minipage} \\
%% \midrule\noalign{}
%% \endhead
%% \bottomrule\noalign{}
%% \endlastfoot
%% Akmola & 57 696 793 & 2 698 159 & 86 183 319 & 540 558 392,1 & 120 723
%% 275 & 22,33 \\
%% Aktobe & 71 801 508 & 2 524 931 & 124 880 322 & 738 017 212,6 & 161 742
%% 664 & 21,92 \\
%% Almaty & 85 030 922 & 6 953 927 & 259 907 517 & 581 733 120,0 & 273 920
%% 111 & 47,09 \\
%% Atyrau & 176 724 848 & 2 105 293 & 510 816 901 & 1 513 836 076,4 & 536
%% 401 972 & 35,43 \\
%% East Kazakhstan & 63 550 133 & 3 130 117 & 154 404 721 & 671 679 073,3 &
%% 174 763 555 & 26,02 \\
%% Zhambyl & 46 482 259 & 3 556 980 & 56 586 304 & 560 033 372,9 & 89 649
%% 652 & 16,01 \\
%% West Kazakhstan & 60 446 743 & 1 783 768 & 134 302 619 & 529 979 578,0 &
%% 156 242 344 & 29,48 \\
%% Karaganda & 100 657 785 & 4 464 760 & 208 881 980 & 1 055 964 178,3 &
%% 251 339 931 & 23,80 \\
%% Kyzylorda & 34 223 426 & 2 060 722 & 64 274 062 & 559 412 725,4 & 81 275
%% 991 & 14,53 \\
%% Kostanay & 62 045 659 & 2 429 879 & 110 444 310 & 615 444 480,0 & 141
%% 786 555 & 23,04 \\
%% Mangistau & 95 760 861 & 3 592 700 & 143 378 655 & 1 004 617 916,4 & 199
%% 718 620 & 19,88 \\
%% Pavlodar & 67 868 655 & 2 415 002 & 156 664 009 & 729 197 218,0 & 179
%% 948 463 & 24,68 \\
%% North Kazakhstan & 33 665 852 & 1 785 494 & 63 242 077 & 342 072 696,0 &
%% 79 720 800 & 23,31 \\
%% Turkestan & 59 251 623 & 4 587 864 & 112 122 541 & 804 988 149,2 & 142
%% 325 266 & 17,68 \\
%% Zhetysu & 25 930 316 & 2 624 607 & 199 383 808 & 320 360 450,6 & 168 123
%% 589 & 52,48 \\
%% Abai & 35 508 593 & 1 863 386 & 46 628 382 & 432 616 630,1 & 70 011 846
%% & 16,18 \\
%% Ulytau & 35 167 663 & 903 658 & -34 306 175 & 363 962 348,2 & 12 056 999
%% & 3,31 \\
%% Shymkent city & 55 307 321 & 4 736 210 & 168 298 785 & 574 604 184,0 &
%% 177 852 681 & 30,95 \\
%% Almaty city & 550 653 936 & 14 727 956 & 2 164 042 653 & 2 611 728 011,3
%% & 2 080 211 749 & 79,65 \\
%% Astana city & 274 609 956 & 7 876 752 & 916 571 353 & 1 802 515 527,5 &
%% 924 086 655 & 51,27 \\
%% \multicolumn{7}{@{}>{\raggedright\arraybackslash}p{(\linewidth - 12\tabcolsep) * \real{1.0000} + 12\tabcolsep}@{}}{%
%% \emph{Note: the table was compiled by the authors based on source
%% {[}16{]}}} \\
%% \end{longtable}

{\bfseries Table 4- Tax burden on households (from the household
perspective)}

%% \begin{longtable}[]{@{}
%%   >{\centering\arraybackslash}p{(\linewidth - 14\tabcolsep) * \real{0.1203}}
%%   >{\centering\arraybackslash}p{(\linewidth - 14\tabcolsep) * \real{0.1322}}
%%   >{\centering\arraybackslash}p{(\linewidth - 14\tabcolsep) * \real{0.1470}}
%%   >{\centering\arraybackslash}p{(\linewidth - 14\tabcolsep) * \real{0.1314}}
%%   >{\centering\arraybackslash}p{(\linewidth - 14\tabcolsep) * \real{0.1321}}
%%   >{\centering\arraybackslash}p{(\linewidth - 14\tabcolsep) * \real{0.1173}}
%%   >{\centering\arraybackslash}p{(\linewidth - 14\tabcolsep) * \real{0.1026}}
%%   >{\centering\arraybackslash}p{(\linewidth - 14\tabcolsep) * \real{0.1174}}@{}}
%% \toprule\noalign{}
%% \endhead
%% \bottomrule\noalign{}
%% \endlastfoot
%% {\bfseries Taxpayer status} & {\bfseries Average per capita income,
%% tenge/month per person} & {\bfseries Income tax, tenge/month per person} &
%% {\bfseries Income after tax, tenge/month per person} & {\bfseries Disposable
%% income, tenge/month per person} & {\bfseries Consumption tax, tenge} &
%% {\bfseries Absolute tax burden, thousand tenge} & {\bfseries Relative tax
%% burden, \%} \\
%% {\bfseries Tax resident} & 210 840 & 21 084 & 189 756 & 170 780 & 20 493,6
%% & 41 577,6 & 19,7 \\
%% {\bfseries Tax non-resident} & 210 840 & 42 168 & 168 672 & 151 804 & 18
%% 216,5 & 60~384,5 & 28,6 \\
%% {\bfseries Averade}
%% 
%% {\bfseries taxpayer} & 210 840 & 52 710 & 179 214 & 161 292 & 19~710,1 &
%% 50~981,1 & 24,2 \\
%% \multicolumn{8}{@{}>{\centering\arraybackslash}p{(\linewidth - 14\tabcolsep) * \real{1.0000} + 14\tabcolsep}@{}}{%
%% \emph{Note: the table was compiled by the authors based on source
%% {[}16{]}}} \\
%% \end{longtable}

Calculations show that the tax burden on tax residents and non-residents
differs in the following areas:

- the tax burden on income tax is higher for non-residents, while the
tax burden on indirect taxes is higher for residents;

- overall, the tax burden (both absolute and relative) is higher for
non-residents and lower for residents.

In their average calculations (for the average taxpayer), the authors
used an average estimated personal income tax rate of (20 + 10) / 2 =
15\%.

The calculated higher burden on non-residents reflects a fiscal strategy
that prioritizes immediate revenue collection from mobile labor
resources. This phenomenon is discussed by Benzarti and Carloni (2019),
who emphasize that the benefits of tax structures are rarely distributed
evenly across all demographic groups. In the Kazakhstani context, the
discrepancy between resident and non-resident burdens underscores the
need for a tax policy that better balances fiscal goals with the need to
attract international human capital {[}13{]}.

{\bfseries Table 5 - Key Analytical Indicators of the Household Tax Burden
in Kazakhstan}

%% \begin{longtable}[]{@{}
%%   >{\raggedright\arraybackslash}p{(\linewidth - 2\tabcolsep) * \real{0.5378}}
%%   >{\raggedright\arraybackslash}p{(\linewidth - 2\tabcolsep) * \real{0.4622}}@{}}
%% \toprule\noalign{}
%% \begin{minipage}[b]{\linewidth}\raggedright
%% Indicator
%% \end{minipage} & \begin{minipage}[b]{\linewidth}\raggedright
%% Value / Result
%% \end{minipage} \\
%% \midrule\noalign{}
%% \endhead
%% \bottomrule\noalign{}
%% \endlastfoot
%% Maximum relative tax burden (by region) & Almaty city (79.65\%) \\
%% Minimum relative tax burden (by region) & Ulytau region (3.31\%) \\
%% Share of expenditures on goods and services & 91.8\% -- 92.8\% (Base for
%% indirect taxes) \\
%% Tax burden on non-residents & 28.6\% (vs.19.7\% for residents) \\
%% \multicolumn{2}{@{}>{\raggedright\arraybackslash}p{(\linewidth - 2\tabcolsep) * \real{1.0000} + 2\tabcolsep}@{}}{%
%% \emph{Note: Compiled by the authors based on modeling results}} \\
%% \end{longtable}

The empirical data presented in Table 5 illustrates the profound
structural and regional imbalances within Kazakhstan' s
tax system. The transition from a descriptive approach to analytical
modeling allows for a deeper understanding of the "hidden" mechanisms of
tax regulation:

\emph{Evidence of Regional Tax Polarity:} The staggering difference
between the relative tax burden in Almaty (79.65\%) and the Ulytau
region (3.31\%) reflects a high concentration of tax capacity in urban
centers and a significant deficit in the tax resource base of newly
formed regions. This confirms the hypothesis that uniform national
monitoring tools are insufficient for such a heterogeneous fiscal
landscape.

\emph{Confirmation of Hidden Regressivity:} The fact that household
expenditures on goods and services remain consistently high (91.8\% --
92.8\%) proves that the population acts as the primary de facto payer of
indirect taxes. Since these taxes (VAT and excises) are not adjusted for
income levels, they create a regressive effect where the relative burden
is higher for low-income groups.

\emph{Analysis of Residency Impact:} The calculated burden on
non-residents (28.6\%) versus residents (19.7\%) reveals that the
current tax policy is primarily revenue-driven in its treatment of
foreign labor, which may create barriers to international human capital
mobility.

These findings move beyond standard statistical reporting by quantifying
the actual tax incidence. By linking regional disparities to consumption
structures, the study provides a scientifically grounded basis for the
proposed methodological modernization of the state' s
fiscal monitoring system.

{\bfseries Conclusion.} This study provides a comprehensive analytical
assessment of the tax regulation system within the household sector of
the Republic of Kazakhstan. The results confirm the initial hypothesis
regarding the existence of profound regional and structural imbalances
that remain obscured by aggregated national statistical averages. By
transitioning from a purely descriptive narrative to analytical
modeling, this research has identified several critical issues that
necessitate immediate methodological and policy intervention.

Firstly, the identified colossal disparity in the relative tax burden -
ranging from a peak of 79.65\% in Almaty to a mere 3.31\% in the Ulytau
region - indicates that the current tax framework fails to account for
the actual geography of consumption and the concentration of financial
flows. Such high variability suggests that the application of uniform
tax monitoring tools across the entire national territory is inefficient
and fails to capture the true fiscal capacity of diverse regions. This
research demonstrates that establishing a fair system for distributing
the fiscal burden is impossible without considering the regional
specificity of the resource base and the environmental "footprint" of
industrial zones.

Secondly, the analysis of household expenditure structures confirmed
that Kazakhstani households allocate up to 92.8\% of their income toward
goods and services. This effectively positions the population as the
primary bearer of indirect taxes, such as VAT and excises. In the
absence of a progressive income tax scale, this creates a "hidden
regressivity" effect, where low-income groups bear a proportionally
higher burden through their daily consumption. This finding aligns with
the global discourse provided by Bachas et al. (2024) and Li Hong
(2019), who emphasize the vulnerability of households in
resource-dependent economies to consumption-based taxation.

Thirdly, the differentiation of the tax burden by residency status -
28.6\% for non-residents versus 19.7\% for residents - demonstrates that
fiscal policy regarding foreign citizens remains primarily
revenue-driven rather than regulatory. This disparity creates additional
barriers to attracting high-quality international human capital and
necessitates a revision of tax deduction mechanisms to foster a more
competitive labor market.

Based on these findings, we recommend that state planning systems
integrate a "full tax incidence" accounting methodology. This approach
must move beyond direct payments to include the calculated share of
indirect taxes relative to the consumer baskets of various social
groups. Furthermore, the digitalization of monitoring through electronic
invoicing should be leveraged to form a detailed database of household
behavior, allowing for more accurate, adaptive, and socially oriented
tax regulation in Kazakhstan.

{\bfseries Литература}

1. Musgrave R.A., Musgrave P.B. Public Finance in Theory and Practice.
-5th ed./ McGraw-Hill. -1989. -664 p. ISBN 978-0070441279, 0070441278

2. Myles G.D. Tax by Design: The Mirrlees Review. -1st ed./ Oxford
University Press. -2011. -552p. ISBN 978-0-19-955374-7, 0199553742.

3. Becker G.S. Accounting for Tastes/Harvard University Press. -1998.
-292p. ISBN 978-0674543577, 0674543572.

4. Piketty T. Capital in the Twenty-First Century/ Harvard University
Press. -2014. -704 p. ISBN 978- 0- 674- 43000- 6, 067443000X.

5. European Commission. Directorate-General for Taxation and Customs
Union. Taxation trends in the European Union: data for the EU Member
States, Iceland, Norway and United Kingdom. -2021 edition/
Publications Office of the European Union. -2021. -274 p. ISBN
978-92-76-32071-5.

6. OECD. Revenue Statistics. Tax Revenue Buoyancy in OECD Countries /OECD
Publishing. -2023.
\href{https://doi.org/10.1787/9d0453d5-en}{DOI:10.1787/9d0453d5-en}.

7. IRS. SOI tax stats: 2023 Statistics of Income Joint Statistical
Research Program projects/Internal Revenue Service. - 2023. URL:
\href{https://www.irs.gov/statistics/soi-tax-stats-2023-statistics-of-income-joint-statistical-research-program-projects.\%20}{https://www.irs.gov/statistics/soi-tax-stats-2023-statistics-of-income-joint-statistical-research-program-projects.}
-Дата обращения: 13.12.2025.

8. Досманбетова~М.С., Халитова~М.М. Қазақстан Республикасының салық
жүктемесінің тиімділігін бағалау // Economy: strategy and practice.
-2020. -15(4). -Б.116-131. DOI:
\href{https://doi.org/10.51176/JESP/issue_4_T11}{10.51176/JESP/issue\_4\_T11}.

9. Нурумов А.А. Система налогообложения Республики Казахстан //
Экономика. Налоги. Право. - 2013. - №4. - С.112-117.

10. Жолмурзаева Г.А., Кожахметова М.К., Кенжебекова Г.Б. Налоговое
администрирование в реализации фискальнои функции
государства//Статистика, учет и аудит. -2020. -№ 3(78). -С.234-239.
ISSN: 1563-2415.

11. Bachas P., Gadenne L., Jensen A. Informality, Consumption Taxes, and
Redistribution //The Review of Economic Studies. -2024. --Vol.91(5).
-Р.2604--2634.
\href{https://doi.org/10.1093/restud/rdad095}{DOI:10.1093/restud/rdad095}.

12. Xiong Zh., Li H. Ecological Deficit Tax: A Tax Design and Simulation
of Compensation for Ecosystem Service Value Based on Ecological
Footprint in China // Journal of Cleaner Production. -2019. -Vol.230.
-P.1128-1137.
\href{https://doi.org/10.1016/j.jclepro.2019.05.172}{DOI:10.1016/j.jclepro.2019.05.172}.

13. Benzarti Y., Carloni D.~Who Really Benefits from Consumption Tax Cuts?
Evidence from a Large VAT Reform in France // American Economic
Journal: Economic Policy. -2019. -Vol.11(1). -P.38-63. DOI:
10.1257/pol.20170504.

14. OECD. Consumption Tax Trends 2024: VAT/GST and Excise, Core Design
Features and Trends/ OECD Publishing. -2024.
\href{https://doi.org/10.1787/dcd4dd36-en}{DOI: 10.1787/dcd4dd36-en}.

15. \href{https://www.gov.kz/memleket/entities/minfin}{Министерство
финансов Республики Казахстан}. Статистический бюллетень на 1 января
2025 года. -URL:
\href{https://www.gov.kz/memleket/entities/minfin/documents/details/792959?lang=ru\%20}{https://www.gov.kz/memleket/entities/minfin/documents/details/792959?lang=ru}.
-Дата обращения: 13.12.2025.

16. Комитет государственных доходов. Фактические поступления по налогам и
платежам в государственный бюджет за 2002 - 2025. URL:
\url{https://kgd.gov.kz/ru/content/fakticheskie-postupleniya-po-nalogam-i-platezham-v-gosudarstvennyy-byudzhet-za-2002-2025-gg}.
-Дата обращения: 13.12.2025.

17. Бюро национальной статистики. Агентства по стратегическому
планированию и реформам Республики Казахстан. Социально-экономическое
развитие Республики Казахстана. Астана. -2025. URL:
\url{https://stat.gov.kz/upload/iblock/5bf/spq44e4kapuxig8n72q1artx5iepxdtv/Ж-01-М\%2001\%202025\%20(рус).pdf?sphrase_id=1656683}.
-Дата обращения: 13.12.2025.

18. Бюро национальной статистики Агентства по стратегическому планированию
и реформам Республики Казахстан. Статистика уровня жизни. Расходы и
доходы населения Республики Казахстан (2024г.) Астана. -2025. URL:
\url{https://stat.gov.kz/ru/industries/labor-and-income/stat-life/publications/333428/}.
-Дата обращения: 13.12.2025.

{\bfseries References}

1. Musgrave R.A., Musgrave P.B. Public Finance in Theory and Practice.
-5th ed./ McGraw-Hill. -1989. -664 p. ISBN 978-0070441279, 0070441278

2. Myles G.D. Tax by Design: The Mirrlees Review. -1st ed./ Oxford
University Press. -2011. -552p. ISBN 978-0-19-955374-7, 0199553742.

3. Becker G.S. Accounting for Tastes/Harvard University Press. -1998.
-292p. ISBN 978-0674543577, 0674543572.

4. Piketty T. Capital in the Twenty-First Century/ Harvard University
Press. -2014. -704 p. ISBN 978- 0- 674- 43000- 6, 067443000X.

5. European Commission. Directorate-General for Taxation and Customs
Union. Taxation trends in the European Union: data for the EU Member
States, Iceland, Norway and United Kingdom. -2021 edition/ Publications
Office of the European Union. -2021. -274 p. ISBN 978-92-76-32071-5.

6. OECD. Revenue Statistics. Tax Revenue Buoyancy in OECD Countries
/OECD Publishing. -2023.
\href{https://doi.org/10.1787/9d0453d5-en}{DOI:10.1787/9d0453d5-en}.

7. IRS. SOI tax stats: 2023 Statistics of Income Joint Statistical
Research Program projects/Internal Revenue Service. - 2023. URL:
\href{https://www.irs.gov/statistics/soi-tax-stats-2023-statistics-of-income-joint-statistical-research-program-projects.\%20}{https://www.irs.gov/statistics/soi-tax-stats-2023-statistics-of-income-joint-statistical-research-program-projects.}
-Дата обращения: 13.12.2025.

8. Dosmanbetova M.S., Halitova M.M. Қazaқstan Respublikasynyң salyқ
zhүktemesіnің tiіmdіlіgіn baғalau // Economy: strategy and practice.
-2020. -15(4). -B.116-131. DOI: 10.51176/JESP/issue\_4\_T11. {[}in
Kazakh{]}

9. Nurumov A.A. Sistema nalogooblozhenija Respubliki Kazahstan //
Jekonomika. Nalogi. Pravo. - 2013. - №4. - S.112-117. {[}in Russian{]}

10. Zholmurzaeva G.A., Kozhahmetova M.K., Kenzhebekova G.B. Nalogovoe
administrirovanie v realizacii fiskal' noi funkcii
gosudarstva//Statistika, uchet i audit. -2020. -№ 3(78). -S.234-239.
ISSN: 1563-2415. {[}in Russian{]}

11. Bachas P., Gadenne L., Jensen A. Informality, Consumption Taxes, and
Redistribution //The Review of Economic Studies. -2024. --Vol.91(5).
-Р.2604--2634.
\href{https://doi.org/10.1093/restud/rdad095}{DOI:10.1093/restud/rdad095}.

12. Xiong Zh., Li H. Ecological Deficit Tax: A Tax Design and Simulation
of Compensation for Ecosystem Service Value Based on Ecological
Footprint in China // Journal of Cleaner Production. -2019. -Vol.230.
-P.1128-1137.
\href{https://doi.org/10.1016/j.jclepro.2019.05.172}{DOI:10.1016/j.jclepro.2019.05.172}.

13. Benzarti Y., Carloni D.~Who Really Benefits from Consumption Tax
Cuts? Evidence from a Large VAT Reform in France // American Economic
Journal: Economic Policy. -2019. -Vol.11(1). -P.38-63. DOI:
10.1257/pol.20170504.

14. OECD. Consumption Tax Trends 2024: VAT/GST and Excise, Core Design
Features and Trends/ OECD Publishing. -2024.
\href{https://doi.org/10.1787/dcd4dd36-en}{DOI: 10.1787/dcd4dd36-en}.

15. Ministerstvo finansov Respubliki Kazahstan. Statisticheskij
bjulleten'{} na 1 janvarja 2025 goda. -URL:
https://www.gov.kz/memleket/entities/minfin/documents/details/792959?lang=ru.
-Data obrashhenija: 13.12.2025. {[}in Russian{]}

16. Komitet gosudarstvennyh dohodov. Fakticheskie postuplenija po nalogam
i platezham v gosudarstvennyj bjudzhet za 2002 - 2025. URL:
https://kgd.gov.kz/ru/content/fakticheskie-postupleniya-po-nalogam-i-platezham-v-gosudarstvennyy-byudzhet-za-2002-2025-gg.
-Data obrashhenija: 13.12.2025. {[}in Russian{]}

17. Bjuro nacional' noj statistiki. Agentstva po
strategicheskomu planirovaniju i reformam Respubliki Kazahstan.
Social' no-jekonomicheskoe razvitie Respubliki
Kazahstana. Astana. -2025. URL:
https://stat.gov.kz/upload/iblock/5bf/spq44e4kapuxig8n72q1artx5iepxdtv/\%D0\%96-01-\%D0\%9C\%2001\%202025\%20(\%D1\%80\%D1\%83\%D1\%81).pdf?sphrase\_id=1656683.
-Data obrashhenija: 13.12.2025. {[}in Russian{]}

18. Bjuro nacional' noj statistiki Agentstva po
strategicheskomu planirovaniju i reformam Respubliki Kazahstan.
Statistika urovnja zhizni. Rashody i dohody naselenija Respubliki
Kazahstan (2024g.) Astana. -2025. URL:
https://stat.gov.kz/ru/industries/labor-and-income/stat-life/publications/333428/.
-Data obrashhenija: 13.12.2025. {[}in Russian{]}

\emph{{\bfseries Information about authors}}

Maulenberdieva G.A. - PhD student, senior lecturer, M. Auezov South
Kazakhstan Research University, Shymkent, Kazakhstan, e-mail:
\href{file:///C:/Users/HP/Desktop/статьи\%20мои/guldar2018g@mail.ru}{guldar2018g@mail.ru};

Abishоva A.U. - candidate of Economic Sciences, Associate Professor, M.
Auezov South Kazakhstan Research University, Shymkent, Kazakhstan,
\href{file:///C:/Users/admin/Downloads/altuka07@mail.ru}{altuka07@mail.ru};

Balabekova D.B - candidate of Economic Sciences, Senior Lecturer,
Central Asian Innovation University, Shymkent, Kazakhstan,
\href{file:///C:/Users/HP/Desktop/статьи\%20мои/ageka1947@mail.ru}{ageka1947@mail.ru};

Zhadigerova G.A. - candidate of Economic Sciences, Senior Lecturer,
Central Asian Innovation University, Shymkent,
Kazakhstan,~\href{file:///C:/Users/HP/Desktop/статьи\%20мои/zhadigerova1020@mail.ru}{zhadigerova1020@mail.ru};

Tuleeva A.Zh. - senior lecturer, M. Auezov South Kazakhstan
Research University, Shymkent, Kazakhstan,
\href{file:///C:/Users/HP/Desktop/статьи\%20мои/nura-meir@mail.ru}{nura-meir@mail.ru}.

~~~~~~

\emph{{\bfseries Информация об авторах}}

Мауленбердиева Г.А. -~докторант, старший преподаватель,
Южно-Казахстанский исследовательский университет им. М. Ауезова,
Шымкент, Казахстан,~e-mail:
\href{file:///C:/Users/HP/Desktop/статьи\%20мои/guldar2018g@mail.ru}{guldar2018g@mail.ru};

Абишова А.У. - кандидат экономических наук, доцент, Южно-Казахстанский
исследовательский университет им. М. Ауезова, Шымкент,
Казахстан,~e-mail:
\href{file:///C:/Users/admin/Downloads/altuka07@mail.ru}{altuka07@mail.ru};

Балабекова Д.Б. - кандидат экономических наук, старший преподаватель,
Центрально-Азиатский инновационный университет, Шымкент,
Казахстан,~e-mail:
\href{file:///C:/Users/HP/Desktop/статьи\%20мои/ageka1947@mail.ru}{ageka1947@mail.ru};

Жадигерова Г.А.~ - кандидат экономических наук, старший преподаватель,
Центрально-Азиатский инновационный университет, Шымкент,
Казахстан,~e-mail:\href{file:///C:/Users/HP/Desktop/статьи\%20мои/zhadigerova1020@mail.ru}{zhadigerova1020@mail.ru};

Тулеева А.Ж. - старший преподаватель, Южно-Казахстанский
исследовательский университет им. М. Ауезова, Шымкент,
Казахстан,~e-mail:
\href{file:///C:/Users/HP/Desktop/статьи\%20мои/nura-meir@mail.ru}{nura-meir@mail.ru}.\